\chapter{Odwzorowania ciągłe}
\label{ch:topog-odwzorowania-ciagle}

\noindent\textbf{Cel rozdziału.}
W rozdziałach poprzednich opisaliśmy przestrzeń przez zbiory otwarte
oraz przez operacje wnętrza, domknięcia i brzegu.
Teraz zbadamy odwzorowania zachowujące tę strukturę.
Wskażemy równoważne kryteria ciągłości, ich postać lokalną i metryczną,
wyprowadzimy kryteria za pomocą baz i podbaz oraz wykażemy
własności składania. Następnie rozróżnimy odwzorowania ciągłe,
otwarte, domknięte, homeomorfizmy i zanurzenia topologiczne.
Szczególnie ważne będą kontrprzykłady: ciągła bijekcja nie musi mieć
ciągłej odwrotności, a ciągłe odwzorowanie nie musi być ani otwarte,
ani domknięte. Nie zakładamy metryzowalności ani żadnych aksjomatów
oddzielania, chyba że wyraźnie je wymieniono.

\section{Ciągłość jako zachowanie zbiorów otwartych}
\label{sec:topog-ciaglosc-definicja}

\begin{definicja}[Odwzorowanie ciągłe]
\label{def:topog-odwzorowanie-ciagle}
Niech $(X,\tau_X)$ i $(Y,\tau_Y)$ będą przestrzeniami topologicznymi.
Odwzorowanie $f\colon X\to Y$ nazywamy \emph{ciągłym}, jeżeli dla
każdego zbioru otwartego $V\subseteq Y$ przeciwobraz
$f^{-1}(V)=\{x\in X:f(x)\in V\}$ jest otwarty w $X$.
Symbol $f^{-1}(V)$ oznacza tu przeciwobraz zbioru, nawet jeśli $f$
nie jest bijekcją. W definicji dopuszczamy $X=\varnothing$
i $Y=\varnothing$, o ile istnieje rozważane odwzorowanie.
\end{definicja}

\begin{uwaga}[Dlaczego przeciwobrazy, a nie obrazy?]
Dla dowolnego odwzorowania $f\colon X\to Y$ i rodziny
$(V_i)_{i\in I}$ podzbiorów $Y$ zachodzą równości
\begin{align*}
f^{-1}\Bigl(\bigcup_{i\in I}V_i\Bigr)
 &=\bigcup_{i\in I}f^{-1}(V_i), &
f^{-1}\Bigl(\bigcap_{i\in I}V_i\Bigr)
 &=\bigcap_{i\in I}f^{-1}(V_i).
\end{align*}
Ponadto $f^{-1}(\varnothing)=\varnothing$ i $f^{-1}(Y)=X$.
Prawa te wynikają bezpośrednio z definicji przeciwobrazu i działają
również dla rodzin pustych. Dzięki temu warunek ciągłości jest
zgodny z aksjomatami topologii. Dla obrazów ogólnego odwzorowania
nie zachodzi analogiczna równość dla przecięć.
\end{uwaga}

\begin{przyklad}[Odwzorowania stałe, identycznościowe i zbiory puste]
\label{ex:topog-ciaglosc-elementarna}
Dla każdej przestrzeni topologicznej $X$ odwzorowanie
$\operatorname{id}_X\colon X\to X$, $\operatorname{id}_X(x)=x$,
jest ciągłe: przeciwobrazem zbioru otwartego jest ten sam zbiór.
Jeżeli $Y$ jest przestrzenią topologiczną, $y_0\in Y$, to
$f\colon X\to Y$, $f(x)=y_0$, jest ciągłe, bo dla dowolnego
$V\subseteq Y$ jego przeciwobrazem jest $X$ albo $\varnothing$.
Każde odwzorowanie ze zbioru pustego do $Y$ jest ciągłe,
ponieważ wszystkie przeciwobrazy są puste. Odwzorowanie
$X\to\varnothing$ istnieje tylko wtedy, gdy $X=\varnothing$.
\end{przyklad}

\begin{przyklad}[Zależność ciągłości od topologii]
\label{ex:topog-ciaglosc-skrajne}
Jeżeli $X$ ma topologię dyskretną, to każde odwzorowanie
$f\colon X\to Y$ jest ciągłe dla dowolnej przestrzeni $Y$:
wszystkie podzbiory $X$ są otwarte.
Jeżeli $Y$ ma topologię antydyskretną, to każde
$f\colon X\to Y$ jest ciągłe, ponieważ jedyne zbiory otwarte
w $Y$ to $\varnothing$ i $Y$.

Niech $S=\{0,1\}$. Odwzorowanie identycznościowe
$\operatorname{id}\colon(S,\mathcal P(S))\to
(S,\{\varnothing,S\})$ jest ciągłą bijekcją. Odwrotne
odwzorowanie nie jest ciągłe: przeciwobraz otwartego zbioru
$\{0\}$ przestrzeni dyskretnej nie jest otwarty w przestrzeni
antydyskretnej.
\end{przyklad}

\begin{twierdzenie}[Równoważne globalne kryteria ciągłości]
\label{thm:topog-ciaglosc-kryteria-globalne}
Niech $X,Y$ będą przestrzeniami topologicznymi, a
$f\colon X\to Y$ odwzorowaniem. Następujące warunki są
równoważne:
\begin{enumerate}[label=(\roman*)]
\item $f$ jest ciągłe;
\item przeciwobraz każdego zbioru domkniętego w $Y$
      jest domknięty w $X$;
\item dla każdego $A\subseteq X$ zachodzi
      $f(\overline A^{\,X})\subseteq\overline{f(A)}^{\,Y}$;
\item dla każdego $B\subseteq Y$ zachodzi
      $\overline{f^{-1}(B)}^{\,X}\subseteq
        f^{-1}(\overline B^{\,Y})$.
\end{enumerate}
\end{twierdzenie}
\begin{proof}
Załóżmy (i). Jeżeli $F\subseteq Y$ jest domknięty, to
$Y\setminus F$ jest otwarty. Z równości
$X\setminus f^{-1}(F)=f^{-1}(Y\setminus F)$ wynika,
że $f^{-1}(F)$ jest domknięty. Otrzymujemy (ii).

Załóżmy (ii), ustalmy $A\subseteq X$ i oznaczmy
$F=\overline{f(A)}^{\,Y}$. Zbiór $F$ jest domknięty,
więc $f^{-1}(F)$ jest domknięty w $X$. Ponieważ
$A\subseteq f^{-1}(F)$, dostajemy
$\overline A^{\,X}\subseteq f^{-1}(F)$, co daje (iii).

Załóżmy (iii) i wybierzmy $B\subseteq Y$.
Stosując (iii) do $A=f^{-1}(B)$, otrzymujemy
\[
f(\overline{f^{-1}(B)}^{\,X})
 \subseteq\overline{f(f^{-1}(B))}^{\,Y}
 \subseteq\overline B^{\,Y}.
\]
Przeciwobraz ostatniego zawierania daje (iv).

Załóżmy (iv) i niech $F\subseteq Y$ będzie domknięty.
Ponieważ $\overline F^{\,Y}=F$, mamy
$\overline{f^{-1}(F)}^{\,X}\subseteq f^{-1}(F)$.
Drugie zawieranie wynika z ekstensywności domknięcia,
więc $f^{-1}(F)$ jest domknięty. Otrzymujemy (ii),
a z (ii) wynika (i), ponieważ dla otwartego $V\subseteq Y$
zbiór $f^{-1}(V)$ jest dopełnieniem domkniętego
$f^{-1}(Y\setminus V)$. Dowodzi to równoważności.
\end{proof}

\begin{stwierdzenie}[Ciągłość i wnętrze]
\label{prop:topog-ciaglosc-wnetrze}
Jeżeli $f\colon X\to Y$ jest ciągłe, to dla każdego
$B\subseteq Y$ zachodzi
$f^{-1}(\operatorname{int}_Y B)
 \subseteq\operatorname{int}_X(f^{-1}(B))$.
\end{stwierdzenie}
\begin{proof}
Zbiór $\operatorname{int}_Y B$ jest otwarty i zawarty w $B$.
Z ciągłości jego przeciwobraz jest otwarty i zawarty w
$f^{-1}(B)$. Ponieważ wnętrze jest największym otwartym
podzbiorem, dostajemy tezę.
\end{proof}

\section{Ciągłość w punkcie i otoczenia}
\label{sec:topog-ciaglosc-lokalna}

\begin{definicja}[Ciągłość w punkcie]
\label{def:topog-ciaglosc-w-punkcie}
Niech $f\colon X\to Y$ będzie odwzorowaniem przestrzeni
topologicznych, a $x\in X$. Mówimy, że $f$ jest
\emph{ciągłe w punkcie $x$}, jeśli dla każdego otwartego
zbioru $V\subseteq Y$ zawierającego $f(x)$ istnieje
otwarty zbiór $U\subseteq X$ taki, że
$x\in U$ i $f(U)\subseteq V$.
\end{definicja}

\begin{twierdzenie}[Kryterium punktowe i otoczeniowe]
\label{thm:topog-ciaglosc-punktowa-globalna}
Dla odwzorowania $f\colon X\to Y$ przestrzeni topologicznych
następujące warunki są równoważne:
\begin{enumerate}[label=(\roman*)]
\item $f$ jest ciągłe;
\item $f$ jest ciągłe w każdym punkcie $x\in X$;
\item dla każdego $x\in X$ i każdego otoczenia $N$
      punktu $f(x)$ zbiór $f^{-1}(N)$ jest otoczeniem $x$.
\end{enumerate}
\end{twierdzenie}
\begin{proof}
Załóżmy (i), ustalmy $x\in X$ i otwarty $V\ni f(x)$.
Przeciwobraz $U=f^{-1}(V)$ jest otwarty, zawiera $x$
i spełnia $f(U)\subseteq V$. Mamy (ii).

Załóżmy (ii). Niech $N$ będzie otoczeniem $f(x)$.
Z definicji otoczenia istnieje otwarty zbiór $V$ taki,
że $f(x)\in V\subseteq N$. Ciągłość w $x$ daje otwarty
$U\ni x$ z $f(U)\subseteq V$. Wtedy
$x\in U\subseteq f^{-1}(N)$, a zatem $f^{-1}(N)$
jest otoczeniem $x$. Otrzymujemy (iii).

Załóżmy (iii) i niech $V\subseteq Y$ będzie otwarty.
Dla każdego $x\in f^{-1}(V)$ zbiór $V$ jest otoczeniem
$f(x)$, więc $f^{-1}(V)$ jest otoczeniem $x$.
Z kryterium otwartości przez otoczenia wynika,
że $f^{-1}(V)$ jest otwarty. Jeśli $f^{-1}(V)=\varnothing$,
jest otwarty na mocy aksjomatów topologii.
Zatem zachodzi (i).
\end{proof}

\begin{twierdzenie}[Metryczne kryterium $\varepsilon$--$\delta$]
\label{thm:topog-ciaglosc-eps-delta}
Niech $(X,d_X)$ i $(Y,d_Y)$ będą przestrzeniami
metrycznymi, wyposażonymi w topologie wyznaczone przez metryki.
Dla $f\colon X\to Y$ i $x\in X$ następujące warunki
są równoważne:
\begin{enumerate}[label=(\roman*)]
\item $f$ jest ciągłe w $x$ w sensie topologicznym;
\item dla każdego $\varepsilon>0$ istnieje $\delta>0$
      takie, że dla wszystkich $z\in X$ z
      $d_X(z,x)<\delta$ zachodzi
      $d_Y(f(z),f(x))<\varepsilon$.
\end{enumerate}
\end{twierdzenie}
\begin{proof}
Załóżmy (i) i ustalmy $\varepsilon>0$.
Kula $B_{d_Y}(f(x),\varepsilon)$ jest otwarta,
więc ciągłość w $x$ daje otwarty $U\ni x$ taki, że
$f(U)\subseteq B_{d_Y}(f(x),\varepsilon)$.
Ponieważ kule metryczne tworzą bazę topologii $X$,
istnieje $\delta>0$ spełniające $B_{d_X}(x,\delta)\subseteq U$.
To jest dokładnie (ii).

Załóżmy (ii). Jeśli $V\subseteq Y$ jest otwarty i
$f(x)\in V$, to istnieje $\varepsilon>0$ takie, że
$B_{d_Y}(f(x),\varepsilon)\subseteq V$.
Dla wybranego z (ii) $\delta>0$ mamy
$f(B_{d_X}(x,\delta))\subseteq V$, więc $f$
jest ciągłe w $x$ w sensie topologicznym.
\end{proof}

\begin{przyklad}[Moduł na prostej]
\label{ex:topog-ciaglosc-modul}
Odwzorowanie $f\colon\R\to\R$, $f(x)=|x|$,
jest ciągłe dla topologii euklidesowych. Dla dowolnych
$x,z\in\R$ z nierówności trójkąta wynikają
$|z|\le|z-x|+|x|$ oraz $|x|\le|x-z|+|z|$.
Zatem $\bigl||z|-|x|\bigr|\le|z-x|$.
Dla danego $\varepsilon>0$ wystarczy przyjąć
$\delta=\varepsilon$ w kryterium $\varepsilon$--$\delta$.
\end{przyklad}

\begin{przyklad}[Funkcja nieciągła]
\label{ex:topog-funkcja-nieciagla}
Na $\R$ z topologią euklidesową rozważmy
$f\colon\R\to\R$ dane wzorem
$f(x)=0$ dla $x\le0$ i $f(x)=1$ dla $x>0$.
Zbiór $V=(1/2,3/2)$ jest otwarty w $\R$, natomiast
$f^{-1}(V)=(0,\infty)$ jest otwarty; to samo
nie wystarcza jeszcze do wykazania ciągłości.
Dla otwartego $W=(-1/2,1/2)$ mamy
$f^{-1}(W)=(-\infty,0]$, który nie jest otwarty.
Zatem $f$ nie jest ciągłe. Bezpośrednio w punkcie
$0$ każde otwarte otoczenie zawiera dodatnie liczby,
na których $f$ ma wartość $1$, więc nie można zawrzeć
jego obrazu w $W$.
\end{przyklad}

\section{Ciągłość sprawdzana na bazach i podbazach}
\label{sec:topog-ciaglosc-bazy}

\begin{twierdzenie}[Kryterium bazy i podbazy w przeciwdziedzinie]
\label{thm:topog-ciaglosc-bazy-podbazy}
Niech $f\colon X\to Y$ będzie odwzorowaniem przestrzeni
topologicznych.
\begin{enumerate}[label=(\roman*)]
\item Jeżeli $\mathcal B$ jest bazą topologii na $Y$,
      to $f$ jest ciągłe wtedy i tylko wtedy, gdy
      $f^{-1}(B)$ jest otwarty w $X$ dla każdego
      $B\in\mathcal B$.
\item Jeżeli $\mathcal S$ jest podbazą topologii na $Y$,
      to $f$ jest ciągłe wtedy i tylko wtedy, gdy
      $f^{-1}(S)$ jest otwarty w $X$ dla każdego
      $S\in\mathcal S$.
\end{enumerate}
\end{twierdzenie}
\begin{proof}
Jeśli $f$ jest ciągłe, obie tezy są oczywiste,
gdyż elementy bazy i podbazy są otwarte.

Załóżmy, że każdy $f^{-1}(B)$ dla $B\in\mathcal B$
jest otwarty. Każdy otwarty $V\subseteq Y$ jest sumą
wszystkich elementów bazy zawartych w $V$. Zatem
\[
f^{-1}(V)=\bigcup_{\substack{B\in\mathcal B\\B\subseteq V}}
 f^{-1}(B),
\]
a więc jest otwarty. Dotyczy to również $V=\varnothing$,
niezależnie od tego, czy $\varnothing\in\mathcal B$.
Dowodzi to (i).

Jeśli przeciwobrazy wszystkich $S\in\mathcal S$
są otwarte, to przeciwobraz każdego skończonego
przecięcia elementów $\mathcal S$ jest otwarty.
Przecięcie pustej rodziny to $Y$, którego przeciwobrazem
jest $X$. Takie skończone przecięcia tworzą bazę
z definicji topologii generowanej przez podbazę.
Z punktu (i) wynika ciągłość, co kończy dowód (ii).
\end{proof}

\begin{przyklad}[Topologia prostej Sorgenfreya]
\label{ex:topog-ciaglosc-sorgenfrey}
Niech $\tau_e$ oznacza topologię euklidesową na $\R$,
a $\tau_{\mathrm{So}}$ topologię generowaną przez
przedziały $[a,b)$, $a<b$. Z rozdziału~1 wiemy, że
$\tau_e\subsetneq\tau_{\mathrm{So}}$.
Odwzorowanie identycznościowe
$\operatorname{id}\colon(\R,\tau_{\mathrm{So}})
\to(\R,\tau_e)$ jest ciągłe, ponieważ każdy zbiór
otwarty w $\tau_e$ jest otwarty także w $\tau_{\mathrm{So}}$.
Odwzorowanie identycznościowe w przeciwnym kierunku
nie jest ciągłe: przeciwobraz $[0,1)$ byłby tym samym
zbiorem, a nie jest on euklidesowo otwarty.
\end{przyklad}

\begin{stwierdzenie}[Porównanie topologii przez identyczność]
\label{prop:topog-porownanie-przez-ciaglosc}
Niech $\tau_1,\tau_2$ będą topologiami na tym samym zbiorze
$X$. Odwzorowanie identycznościowe
$\operatorname{id}\colon(X,\tau_1)\to(X,\tau_2)$
jest ciągłe wtedy i tylko wtedy, gdy $\tau_2\subseteq\tau_1$.
\end{stwierdzenie}
\begin{proof}
Dla każdego $V\subseteq X$ zachodzi
$\operatorname{id}^{-1}(V)=V$. Warunek ciągłości
mówi więc dokładnie, że każdy zbiór należący do
$\tau_2$ należy do $\tau_1$.
\end{proof}

\begin{przyklad}[Przestrzeń Sierpińskiego wykrywa otwartość]
\label{ex:topog-sierpinski-klasyfikator}
Niech $S=\{0,1\}$ ma topologię
$\{\varnothing,\{1\},S\}$, a $A\subseteq X$.
Określmy odwzorowanie $\chi_A\colon X\to S$ wzorem
$\chi_A(x)=1$ dla $x\in A$ i $\chi_A(x)=0$ dla $x\notin A$.
Ponieważ $\chi_A^{-1}(\{1\})=A$,
a przeciwobrazy dwóch pozostałych otwartych zbiorów
są zawsze otwarte, odwzorowanie $\chi_A$ jest ciągłe
dokładnie wtedy, gdy $A$ jest otwarty w $X$.
W szczególności przykład działa także dla $X=\varnothing$.
\end{przyklad}

\section{Składanie, ograniczenia i sklejanie}
\label{sec:topog-ciaglosc-konstrukcje}

\begin{twierdzenie}[Składanie odwzorowań ciągłych]
\label{thm:topog-skladanie-ciaglych}
Niech $X,Y,Z$ będą przestrzeniami topologicznymi, a
$f\colon X\to Y$, $g\colon Y\to Z$ odwzorowaniami
ciągłymi. Wtedy $g\circ f\colon X\to Z$ jest ciągłe.
\end{twierdzenie}
\begin{proof}
Jeżeli $W\subseteq Z$ jest otwarty, to $g^{-1}(W)$
jest otwarty w $Y$, a więc $f^{-1}(g^{-1}(W))$
jest otwarty w $X$. Ponieważ
$(g\circ f)^{-1}(W)=f^{-1}(g^{-1}(W))$,
zachodzi definicja ciągłości złożenia.
\end{proof}

\begin{definicja}[Topologia podprzestrzeni]
\label{def:topog-topologia-podprzestrzeni-wstep}
Jeśli $(X,\tau)$ jest przestrzenią topologiczną,
a $A\subseteq X$, to na $A$ określamy rodzinę
$\tau_A=\{A\cap U:U\in\tau\}$. Jest ona topologią:
zawiera $\varnothing$ i $A$, zaś przecięcia ze zbiorem
$A$ zachowują dowolne sumy i skończone przecięcia.
Nazywamy ją \emph{topologią podprzestrzeni}.
Konstrukcję tę zbadamy szczegółowo w rozdziale~4.
\end{definicja}

\begin{stwierdzenie}[Ograniczanie dziedziny i przeciwdziedziny]
\label{prop:topog-ograniczanie-ciaglosci}
Niech $f\colon X\to Y$ będzie odwzorowaniem ciągłym.
\begin{enumerate}[label=(\roman*)]
\item Dla każdego $A\subseteq X$ ograniczenie
      $f|_A\colon A\to Y$ jest ciągłe, gdy $A$
      ma topologię podprzestrzeni.
\item Jeżeli $B\subseteq Y$ i $f(X)\subseteq B$,
      to odwzorowanie $\widetilde f\colon X\to B$,
      $\widetilde f(x)=f(x)$, jest ciągłe,
      gdy $B$ ma topologię podprzestrzeni.
\end{enumerate}
\end{stwierdzenie}
\begin{proof}
Dla dowolnego otwartego $V\subseteq Y$ mamy
$(f|_A)^{-1}(V)=A\cap f^{-1}(V)$, co jest otwarte
w podprzestrzeni $A$. Dowodzi to (i).

Jeżeli $W\subseteq B$ jest otwarty w $B$,
to ma postać $B\cap V$ z $V$ otwartym w $Y$.
Ponieważ $f(X)\subseteq B$, zachodzi
$\widetilde f^{-1}(W)=f^{-1}(V)$, zbiór otwarty
w $X$. Otrzymujemy (ii).
\end{proof}

\begin{twierdzenie}[Lokalność ciągłości na otwartym pokryciu]
\label{thm:topog-ciaglosc-lokalna-pokrycie}
Niech $(U_i)_{i\in I}$ będzie rodziną otwartych
podzbiorów przestrzeni topologicznej $X$ taką, że
$X=\bigcup_{i\in I}U_i$. Niech $f\colon X\to Y$ będzie
odwzorowaniem do przestrzeni topologicznej $Y$.
Wówczas $f$ jest ciągłe wtedy i tylko wtedy, gdy
dla każdego $i\in I$ ograniczenie
$f|_{U_i}\colon U_i\to Y$ jest ciągłe.
\end{twierdzenie}
\begin{proof}
Jeśli $f$ jest ciągłe, ciągłość wszystkich ograniczeń
wynika z poprzedniego stwierdzenia.
Odwrotnie, niech $V\subseteq Y$ będzie otwarty.
Dla każdego $i$ zbiór
$U_i\cap f^{-1}(V)=(f|_{U_i})^{-1}(V)$ jest otwarty
w $U_i$. Ponieważ $U_i$ jest otwarty w $X$,
każdy zbiór otwarty w $U_i$ jest także otwarty w $X$.
Zatem $f^{-1}(V)=\bigcup_{i\in I}(U_i\cap f^{-1}(V))$
jest otwarty w $X$. Dla $I=\varnothing$ pokrycie
wymusza $X=\varnothing$ i ten sam argument działa.
\end{proof}

\begin{twierdzenie}[Lemat o sklejaniu na pokryciu skończonym domkniętym]
\label{thm:topog-sklejanie-domkniete}
Niech $X$ i $Y$ będą przestrzeniami topologicznymi,
a $F_1,\ldots,F_m$ domkniętymi podzbiorami $X$, gdzie
$m\ge1$ i $X=F_1\cup\cdots\cup F_m$.
Jeżeli $f\colon X\to Y$ jest odwzorowaniem takim, że
$f|_{F_j}\colon F_j\to Y$ jest ciągłe dla każdego
$j=1,\ldots,m$, to $f$ jest ciągłe.
\end{twierdzenie}
\begin{proof}
Niech $C\subseteq Y$ będzie domknięty. Dla każdego
$j$ zbiór $F_j\cap f^{-1}(C)$ jest domknięty
w podprzestrzeni $F_j$. Jego dopełnienie w $F_j$ jest
otwarte w $F_j$, czyli ma postać $F_j\cap U_j$
dla pewnego otwartego $U_j\subseteq X$.
Zatem $F_j\cap f^{-1}(C)=F_j\cap E_j$,
gdzie $E_j=X\setminus U_j$ jest domknięty w $X$.
Ponieważ $F_j$ jest domknięty w $X$, zbiór
$F_j\cap E_j$ jest domknięty w $X$.
Równość $X=\bigcup_{j=1}^{m}F_j$ daje
\[
f^{-1}(C)=\bigcup_{j=1}^{m}(F_j\cap f^{-1}(C)),
\]
więc $f^{-1}(C)$ jest domknięty jako skończona suma
zbiorów domkniętych. Ciągłość wynika z kryterium
przeciwobrazów zbiorów domkniętych.
\end{proof}

\begin{przyklad}[Dlaczego skończoność pokrycia domkniętego jest istotna?]
\label{ex:topog-sklejanie-nieskonczone}
Niech $X=Y=\R$ mają topologię euklidesową i niech
$f(x)=1$ dla $x=0$ oraz $f(x)=0$ dla $x\ne0$.
Funkcja nie jest ciągła, ponieważ
$f^{-1}((1/2,3/2))=\{0\}$ nie jest otwarty.
Tymczasem każdy singleton $F_x=\{x\}$ jest domknięty
w $\R$, rodzina $(F_x)_{x\in\R}$ pokrywa $\R$,
a wszystkie ograniczenia $f|_{F_x}$ są stałe,
więc ciągłe. Nie można więc bez dodatkowych założeń
zastąpić skończonej rodziny domkniętej w lemacie
rodziną dowolną.
\end{przyklad}

\section{Odwzorowania otwarte i domknięte}
\label{sec:topog-otwarte-domkniete-odwzorowania}

\begin{definicja}[Odwzorowanie otwarte i domknięte]
\label{def:topog-otwarte-domkniete-odwzorowanie}
Niech $f\colon X\to Y$ będzie odwzorowaniem przestrzeni
topologicznych. Nazywamy je \emph{otwartym}, jeśli
$f(U)$ jest otwarty w $Y$ dla każdego otwartego
$U\subseteq X$. Nazywamy je \emph{domkniętym},
jeśli $f(F)$ jest domknięty w $Y$ dla każdego
domkniętego $F\subseteq X$.
Żadna z tych definicji nie zawiera założenia ciągłości
ani surjektywności $f$.
\end{definicja}

\begin{stwierdzenie}[Kryteria przez wnętrze i domknięcie]
\label{prop:topog-otwarte-domkniete-kryteria}
Dla dowolnego odwzorowania $f\colon X\to Y$
przestrzeni topologicznych zachodzą równoważności:
\begin{enumerate}[label=(\roman*)]
\item $f$ jest otwarte wtedy i tylko wtedy, gdy
      $f(\operatorname{int}_X A)
        \subseteq\operatorname{int}_Y f(A)$
      dla każdego $A\subseteq X$;
\item $f$ jest domknięte wtedy i tylko wtedy, gdy
      $\overline{f(A)}^{\,Y}\subseteq f(\overline A^{\,X})$
      dla każdego $A\subseteq X$.
\end{enumerate}
\end{stwierdzenie}
\begin{proof}
Jeśli $f$ jest otwarte, to $f(\operatorname{int}A)$
jest otwarty w $Y$ i zawarty w $f(A)$, więc
jest zawarty we wnętrzu $f(A)$. Odwrotnie,
dla otwartego $U\subseteq X$ mamy
$\operatorname{int}_X U=U$, a założenie daje
$f(U)\subseteq\operatorname{int}_Y f(U)$.
Odwrotne zawieranie zawsze zachodzi, zatem
$f(U)$ jest otwarty. Dowodzi to (i).

Jeśli $f$ jest domknięte, to $f(\overline A^{\,X})$
jest domknięty w $Y$ i zawiera $f(A)$, więc
$\overline{f(A)}^{\,Y}\subseteq f(\overline A^{\,X})$.
Odwrotnie, dla dowolnego domkniętego $F\subseteq X$
założenie daje
$\overline{f(F)}^{\,Y}\subseteq f(\overline F^{\,X})=f(F)$.
Druga inkluzja zachodzi zawsze, więc $f(F)$ jest
domknięty. Dowodzi to (ii).
\end{proof}

\begin{przyklad}[Ciągłość nie wymusza otwartości ani domkniętości]
\label{ex:topog-ciaglosc-a-otwartosc-domknietosc}
Niech $S=\{0,1\}$ i
$f\colon(S,\mathcal P(S))\to(S,\{\varnothing,S\})$
będzie odwzorowaniem identycznościowym.
Jak już wykazaliśmy, $f$ jest ciągłe. Nie jest ono
otwarte: zbiór $\{0\}$ jest otwarty w dziedzinie,
a jego obraz nie jest otwarty w przeciwdziedzinie.
Nie jest też domknięte: zbiór $\{0\}$ jest domknięty
w dziedzinie, a jego obraz nie jest domknięty
w topologii antydyskretnej.

W drugą stronę, identyczność
$g\colon(S,\{\varnothing,S\})\to(S,\mathcal P(S))$
jest otwarta i domknięta (obrazami zbiorów otwartych,
odpowiednio domkniętych, są $\varnothing$ i $S$),
ale nie jest ciągła, gdyż przeciwobraz $\{0\}$ nie
jest otwarty w dziedzinie.
\end{przyklad}

\begin{stwierdzenie}[Składanie odwzorowań otwartych i domkniętych]
\label{prop:topog-skladanie-otwartych-domknietych}
Jeżeli $f\colon X\to Y$ i $g\colon Y\to Z$ są
odwzorowaniami otwartymi, to $g\circ f$ jest otwarte.
Jeżeli oba są domknięte, to $g\circ f$ jest domknięte.
\end{stwierdzenie}
\begin{proof}
Dla otwartego $U\subseteq X$ zbiór $f(U)$ jest otwarty
w $Y$, a zatem $g(f(U))$ jest otwarty w $Z$.
Dla domkniętego $F\subseteq X$ analogicznie
$f(F)$ jest domknięty w $Y$, a $g(f(F))$ w $Z$.
\end{proof}

\section{Homeomorfizmy i zanurzenia topologiczne}
\label{sec:topog-homeomorfizmy}

\begin{definicja}[Homeomorfizm]
\label{def:topog-homeomorfizm}
Niech $X,Y$ będą przestrzeniami topologicznymi.
\emph{Homeomorfizmem} $f\colon X\to Y$ nazywamy
odwzorowanie bijektywne, ciągłe i takie, że
odwzorowanie odwrotne $f^{-1}\colon Y\to X$ jest ciągłe.
Przestrzenie, między którymi istnieje homeomorfizm,
nazywamy \emph{homeomorficznymi} i piszemy $X\cong Y$.
\end{definicja}

\begin{twierdzenie}[Ciągła bijekcja: kryteria homeomorfizmu]
\label{thm:topog-homeomorfizm-kryteria}
Niech $f\colon X\to Y$ będzie ciągłą bijekcją
przestrzeni topologicznych. Następujące warunki są
równoważne:
\begin{enumerate}[label=(\roman*)]
\item $f$ jest homeomorfizmem;
\item $f$ jest odwzorowaniem otwartym;
\item $f$ jest odwzorowaniem domkniętym.
\end{enumerate}
\end{twierdzenie}
\begin{proof}
Niech $h=f^{-1}\colon Y\to X$ będzie odwrotnością
w sensie odwzorowania, istniejącą dzięki bijektywności $f$.
Dla dowolnego $U\subseteq X$ zachodzi
$h^{-1}(U)=f(U)$, dlatego $h$ jest ciągłe
dokładnie wtedy, gdy $f$ jest otwarte. To dowodzi
równoważności (i) i (ii).

Podobnie, stosując kryterium przeciwobrazów zbiorów
domkniętych do $h$, otrzymujemy, że $h$ jest ciągłe
dokładnie wtedy, gdy $f(F)$ jest domknięte dla każdego
domkniętego $F\subseteq X$. To jest warunek (iii).
\end{proof}

\begin{przyklad}[Ciągła bijekcja niebędąca homeomorfizmem]
\label{ex:topog-bijekcja-niehomeomorfizm}
Identyczność
$f\colon(S,\mathcal P(S))\to(S,\{\varnothing,S\})$
dla dwuelementowego $S$ jest ciągłą bijekcją,
ale nie jest homeomorfizmem: jej odwrotność nie jest
ciągła. W świetle poprzedniego twierdzenia wynika to
również z tego, że $f$ nie jest otwarte.
Przykład nie wymaga ani metryki, ani nieskończonego zbioru.
\end{przyklad}

\begin{stwierdzenie}[Zachowanie operacji zbiorowych]
\label{prop:topog-homeomorfizm-operacje}
Niech $f\colon X\to Y$ będzie homeomorfizmem przestrzeni
topologicznych. Wówczas dla każdego $A\subseteq X$
zachodzą równości:
\begin{align*}
f(\operatorname{int}_X A)&=\operatorname{int}_Y f(A),\\
f(\overline A^{\,X})&=\overline{f(A)}^{\,Y},\\
f(\partial_X A)&=\partial_Y f(A),\\
f(A')&=(f(A))'.
\end{align*}
W ostatniej równości symbole $A'$ i $(f(A))'$ oznaczają
zbiory pochodne obliczone odpowiednio w $X$ i w $Y$.
W szczególności homeomorfizmy zachowują gęstość,
domkniętość, otwartość i własność bycia punktem izolowanym
podzbioru.
\end{stwierdzenie}
\begin{proof}
Homeomorfizm $f$ i jego odwrotność są otwarte.
Zatem $f(\operatorname{int}_X A)$ jest otwarty
i zawarty w $f(A)$, co daje
$f(\operatorname{int}_X A)\subseteq\operatorname{int}_Y f(A)$.
Tę samą inkluzję dla $f^{-1}$ i $f(A)$
przekształcamy w odwrotną, otrzymując równość.

Ciągłość $f$ daje
$f(\overline A^{\,X})\subseteq\overline{f(A)}^{\,Y}$,
a ciągłość $f^{-1}$ daje inkluzję odwrotną.
Ponieważ $f$ jest bijekcją, mamy
$f(X\setminus A)=Y\setminus f(A)$
i $f(B\cap C)=f(B)\cap f(C)$.
Ze wzoru
$\partial_X A=\overline A^{\,X}\cap
\overline{X\setminus A}^{\,X}$
wynika więc równość dla brzegów.

Dla $x\in X$ warunek $x\in A'$ oznacza,
że każde otwarte otoczenie $U\ni x$
spotyka $A\setminus\{x\}$.
Homeomorfizm przeprowadza otoczenia $x$
bijektywnie na otoczenia $f(x)$, przy czym
\[
f\bigl((U\setminus\{x\})\cap A\bigr)
=(f(U)\setminus\{f(x)\})\cap f(A).
\]
Daje to równoważność
$x\in A'\Longleftrightarrow f(x)\in(f(A))'$.
Zachowanie gęstości, otwartości i domkniętości
wynika z równości dotyczących domknięcia
i wnętrza, natomiast punkt $x\in A$
jest izolowany w $A$ dokładnie wtedy,
gdy $x\notin A'$.
\end{proof}

\begin{definicja}[Zanurzenie topologiczne]
\label{def:topog-zanurzenie-topologiczne}
Niech $X,Y$ będą przestrzeniami topologicznymi.
Odwzorowanie $f\colon X\to Y$ nazywamy
\emph{zanurzeniem topologicznym}, jeżeli odwzorowanie
$f\colon X\to f(X)$ (z topologią podprzestrzeni
odziedziczoną z $Y$) jest homeomorfizmem.
W szczególności zanurzenie topologiczne jest ciągłe
i różnowartościowe, ale nie musi być surjektywne
na całe $Y$.
\end{definicja}

\begin{stwierdzenie}[Kryterium zanurzenia]
\label{prop:topog-zanurzenie-kryterium}
Niech $f\colon X\to Y$ będzie odwzorowaniem
przestrzeni topologicznych. Jest ono zanurzeniem
topologicznym wtedy i tylko wtedy, gdy jest różnowartościowe,
ciągłe i dla każdego otwartego $U\subseteq X$ zbiór
$f(U)$ jest otwarty w podprzestrzeni $f(X)\subseteq Y$.
\end{stwierdzenie}
\begin{proof}
Jeśli $f$ jest zanurzeniem, to jako homeomorfizm
$X\to f(X)$ jest różnowartościowe, ciągłe i otwarte
do $f(X)$, więc ma podane własności.
Odwrotnie, jeżeli $f$ jest różnowartościowe,
to $f\colon X\to f(X)$ jest bijekcją.
Jest ono ciągłe, ponieważ dla otwartego
$W=f(X)\cap V$ w $f(X)$, z $V$ otwartym w $Y$,
mamy $f^{-1}(W)=f^{-1}(V)$, otwarty w $X$.
Z założenia jest także otwarte do $f(X)$.
Kryterium homeomorfizmu dowodzi tezy.
\end{proof}

\begin{przyklad}[Naturalne włożenie i kontrprzykład]
\label{ex:topog-zanurzenie-przyklady}
Włożenie podprzestrzeni $A\hookrightarrow Y$ jest zanurzeniem.
Ciągła różnowartościowa tożsamość z
Przykładu~\ref{ex:topog-bijekcja-niehomeomorfizm}
nie jest nim, gdyż jej odwrotność nie jest ciągła.
\end{przyklad}

\section{Podsumowanie i zależności}
\label{sec:topog-ciaglosc-podsumowanie}
Ciągłość opisują przeciwobrazy zbiorów otwartych
i domkniętych, zachowanie domknięcia, otoczenia,
a w przestrzeniach metrycznych warunek
$\varepsilon$--$\delta$; wystarcza baza lub podbaza.
Zachowują ją składanie, ograniczanie i sklejanie
na pokryciach otwartych lub skończonych domkniętych.
Otwartość i domkniętość nie gwarantują ciągłości.
Homeomorfizm ma ciągłą odwrotność,
zanurzenie jest homeomorfizmem na obraz.
Dalej rozpatrujemy podprzestrzenie, produkty,
sumy rozłączne, topologie początkowe i końcowe,
zwartość, spójność oraz zbieżność.
